\begin{proof}[Proof of \cref{obj:lem:baseline-translation-affinity}]
\begin{enumerate}
\item \emph{The global Lipschitz bound.}
Define the clamping map on \(\mathbb R\) by
\[
\operatorname{clamp}(w)
=\max\{-1/4,\min\{1/4,w\}\}.
\]
On the support interval, \(\cos(2\pi w)\ge0\), so taking the square root of \(f(w)\) gives \(2\cos(2\pi w)\). Below the interval, the clamping map equals \(-1/4\); above it, the clamping map equals \(1/4\). Since the cosine vanishes at both corresponding angles, these cases give the identity
\[
\varphi(w)=2\cos\bigl(2\pi\operatorname{clamp}(w)\bigr)
\qquad(w\in\mathbb R).
\]
Taking a minimum or maximum with a fixed real number preserves the Lipschitz constant \(1\). Consequently, for \(w_1,w_2\in\mathbb R\),
\[
\bigl|\operatorname{clamp}(w_1)-\operatorname{clamp}(w_2)\bigr|
\le |w_1-w_2|.
\]
Using \(|\cos \omega_1-\cos \omega_2|\le|\omega_1-\omega_2|\) for real angles \(\omega_1,\omega_2\), we obtain
\[
\begin{aligned}
|\varphi(w_1)-\varphi(w_2)|
&\le 2\bigl|2\pi\operatorname{clamp}(w_1)
             -2\pi\operatorname{clamp}(w_2)\bigr|\\
&=4\pi\bigl|\operatorname{clamp}(w_1)
             -\operatorname{clamp}(w_2)\bigr|\\
&\le4\pi|w_1-w_2|.
\end{aligned}
\]
% lean: CausalSmith.Experimentation.ThinnedgraphAdditiveRiskFrontier.baseline_sqrt_lipschitz

\item \emph{The sine energy.}
Changing variables to the real angle \(\omega=2\pi w\), and using the antiderivative
\(\omega/2-\sin\omega\cos\omega/2\) of \(\sin^2\omega\), gives
\[
\begin{aligned}
\int_{-1/4}^{1/4}(4\pi\sin(2\pi w))^2\,dw
&=8\pi\int_{-\pi/2}^{\pi/2}\sin^2\omega\,d\omega\\
&=8\pi
  \left[\frac{\omega}{2}
        -\frac{\sin\omega\cos\omega}{2}\right]_{-\pi/2}^{\pi/2}\\
&=4\pi^2.
\end{aligned}
\]
The closed-interval Lebesgue integral has the same value as this interval integral because its endpoints have measure zero.
% lean: CausalSmith.Experimentation.ThinnedgraphAdditiveRiskFrontier.baseline_derivative_energy

\item \emph{An integral representation.}
For the translation argument, define the closed-support function
\[
\psi(w)=
\begin{cases}
-4\pi\sin(2\pi w),& |w|\le1/4,\\
0,& |w|>1/4,
\end{cases}
\qquad w\in\mathbb R.
\]
Both \(\psi\) and \(\psi^2\) are measurable and integrable by compact support and continuity of the sine expression. The preceding calculation gives
\[
\int_{\mathbb R}\psi(w)^2\,dw
=\int_{-1/4}^{1/4}(4\pi\sin(2\pi w))^2\,dw
=4\pi^2.
\]

Fix \(\rho_0,\rho_1\in\mathbb R\) with \(\rho_0\le\rho_1\), and define
\[
\rho_L=\max\{-1/4,\rho_0\},
\qquad
\rho_U=\min\{1/4,\rho_1\}.
\]
Restricting the integral to the support yields
\[
\int_{\rho_0}^{\rho_1}\psi(w)\,dw
=\int_{[\rho_L,\rho_U]}-4\pi\sin(2\pi w)\,dw,
\]
where the set interval is empty when \(\rho_U<\rho_L\).
If \(\rho_1<-1/4\), this interval is empty and both square-root density values are zero. If \(\rho_0>1/4\), the same conclusion holds. In the remaining case,
\(\rho_1\ge-1/4\) and \(\rho_0\le1/4\), so \(\rho_L\le\rho_U\). Integrating the derivative of \(2\cos(2\pi w)\) gives
\[
\begin{aligned}
\int_{\rho_0}^{\rho_1}\psi(w)\,dw
&=2\cos(2\pi\rho_U)-2\cos(2\pi\rho_L)\\
&=\varphi(\rho_1)-\varphi(\rho_0),
\end{aligned}
\]
because \(\rho_U=\operatorname{clamp}(\rho_1)\) and
\(\rho_L=\operatorname{clamp}(\rho_0)\) in this case. Thus, for every ordered pair of real endpoints,
\[
\int_{\rho_0}^{\rho_1}\psi(w)\,dw
=\varphi(\rho_1)-\varphi(\rho_0).
\]
% lean: CausalSmith.Experimentation.ThinnedgraphAdditiveRiskFrontier.baselineSqrtDerivative_integral

\item \emph{The translation estimate.}
Fix \(w\in\mathbb R\) and \(\rho_0\le\rho_1\), and define
\[
\mathcal I_{\mathrm{int}}(w;\rho_0,\rho_1)
=\int_{\rho_0}^{\rho_1}\psi(w+\rho)\,d\rho.
\]
When \(\rho_0=\rho_1\), this integral is zero. When \(\rho_0<\rho_1\), expanding a nonnegative square gives
\[
\begin{aligned}
0
&\le
\int_{\rho_0}^{\rho_1}
\left((\rho_1-\rho_0)\psi(w+\rho)
      -\mathcal I_{\mathrm{int}}(w;\rho_0,\rho_1)\right)^2\,d\rho\\
&=(\rho_1-\rho_0)
\left[
(\rho_1-\rho_0)\int_{\rho_0}^{\rho_1}\psi(w+\rho)^2\,d\rho
-\mathcal I_{\mathrm{int}}(w;\rho_0,\rho_1)^2
\right].
\end{aligned}
\]
Dividing by the positive interval length, and including the zero-length case, proves
\[
\mathcal I_{\mathrm{int}}(w;\rho_0,\rho_1)^2
\le
(\rho_1-\rho_0)\int_{\rho_0}^{\rho_1}\psi(w+\rho)^2\,d\rho.
\]
The integral representation from the preceding step identifies the left side with
\((\varphi(w+\rho_1)-\varphi(w+\rho_0))^2\).
Translation invariance gives
\[
\int_{\mathbb R}\psi(w+\rho)^2\,dw
=\int_{\mathbb R}\psi(w)^2\,dw.
\]
In particular, the square integrand is integrable on
\(\mathbb R\times[\rho_0,\rho_1]\). Integrating the pointwise bound and applying Fubini therefore gives
\[
\begin{aligned}
\int_{\mathbb R}
(\varphi(w+\rho_1)-\varphi(w+\rho_0))^2\,dw
&\le
(\rho_1-\rho_0)
\int_{\mathbb R}\int_{\rho_0}^{\rho_1}
\psi(w+\rho)^2\,d\rho\,dw\\
&=(\rho_1-\rho_0)
\int_{\rho_0}^{\rho_1}\int_{\mathbb R}
\psi(w+\rho)^2\,dw\,d\rho\\
&=(\rho_1-\rho_0)^2\int_{\mathbb R}\psi(w)^2\,dw\\
&=4\pi^2(\rho_1-\rho_0)^2.
\end{aligned}
\]
For real shifts \(\rho_{\mathrm{shift},1},\rho_{\mathrm{shift},2}\) with \(\rho_{\mathrm{shift},2}\le\rho_{\mathrm{shift},1}\), take \(\rho_0=-\rho_{\mathrm{shift},1}\) and \(\rho_1=-\rho_{\mathrm{shift},2}\). Reversing the difference inside the square then gives
\[
\int_{\mathbb R}(\varphi(w-\rho_{\mathrm{shift},1})-\varphi(w-\rho_{\mathrm{shift},2}))^2\,dw
\le4\pi^2(\rho_{\mathrm{shift},1}-\rho_{\mathrm{shift},2})^2.
\]
For \(\rho_{\mathrm{shift},1}\le\rho_{\mathrm{shift},2}\), apply the same argument with the shifts exchanged; both squared differences are unchanged. This covers all \(\rho_{\mathrm{shift},1},\rho_{\mathrm{shift},2}\), including equality.
% lean: CausalSmith.Experimentation.ThinnedgraphAdditiveRiskFrontier.baseline_translation_energy

\item \emph{Products of translated densities.}
Fix \(B\ge0\) and vectors \(\boldsymbol s,\boldsymbol t\in\mathbb R^B\). For \(1\le\ell\le B\) and \(w\in\mathbb R\), define
\[
f_{\ell,\boldsymbol s}(w)=f(w-s_\ell),
\qquad
f_{\ell,\boldsymbol t}(w)=f(w-t_\ell).
\]
The baseline density is nonnegative and integrable. Its normalization follows from
\[
\begin{aligned}
\int_{\mathbb R}f(w)\,dw
&=\frac{2}{\pi}\int_{-\pi/2}^{\pi/2}\cos^2\omega\,d\omega\\
&=\frac{2}{\pi}
  \left[\frac{\omega}{2}
        +\frac{\sin\omega\cos\omega}{2}\right]_{-\pi/2}^{\pi/2}
=1.
\end{aligned}
\]
Translation preserves integrability and the integral, so
\[
\int_{\mathbb R}f_{\ell,\boldsymbol s}(w)\,dw
=\int_{\mathbb R}f_{\ell,\boldsymbol t}(w)\,dw=1.
\]
Define the coordinate affinities and energies by
\[
\alpha_{\mathrm{aff},\ell}
=\int_{\mathbb R}
\sqrt{f_{\ell,\boldsymbol s}(w)f_{\ell,\boldsymbol t}(w)}\,dw,
\qquad
\mathcal E_{\mathrm{Hel},\ell}
=\int_{\mathbb R}
\left(\sqrt{f_{\ell,\boldsymbol s}(w)}
      -\sqrt{f_{\ell,\boldsymbol t}(w)}\right)^2\,dw.
\]
The square roots are square integrable, so their product is integrable. Expanding the square and using normalization yields
\[
\mathcal E_{\mathrm{Hel},\ell}
=2(1-\alpha_{\mathrm{aff},\ell}).
\]
The affinity is nonnegative, and the energy is nonnegative; hence
\[
0\le\alpha_{\mathrm{aff},\ell}\le1.
\]

For \(y\in\mathbb R^B\), define the product densities
\[
v_{\boldsymbol s}(y)=\prod_{\ell=1}^B f_{\ell,\boldsymbol s}(y_\ell),
\qquad
v_{\boldsymbol t}(y)=\prod_{\ell=1}^B f_{\ell,\boldsymbol t}(y_\ell).
\]
Product integration shows that both densities are integrable and normalized:
\[
\int_{\mathbb R^B}v_{\boldsymbol s}(y)\,dy
=\prod_{\ell=1}^B\int_{\mathbb R}f_{\ell,\boldsymbol s}(w)\,dw=1,
\]
and likewise for \(v_{\boldsymbol t}\).
Nonnegativity allows the square root to factor, so another application of product integration gives
\[
\begin{aligned}
\int_{\mathbb R^B}\sqrt{v_{\boldsymbol s}(y)v_{\boldsymbol t}(y)}\,dy
&=\int_{\mathbb R^B}
\prod_{\ell=1}^B
\sqrt{f_{\ell,\boldsymbol s}(y_\ell)
      f_{\ell,\boldsymbol t}(y_\ell)}\,dy\\
&=\prod_{\ell=1}^B\alpha_{\mathrm{aff},\ell}.
\end{aligned}
\]
Expanding the product-density square therefore gives
\[
\int_{\mathbb R^B}
\left(\sqrt{v_{\boldsymbol s}(y)}
      -\sqrt{v_{\boldsymbol t}(y)}\right)^2\,dy
=2\left(1-\prod_{\ell=1}^B\alpha_{\mathrm{aff},\ell}\right).
\]

For completeness, the product-defect bound follows by induction over subsets
\(\mathcal J\subseteq\{1,\ldots,B\}\). Every such product lies in \([0,1]\), and the empty subset gives equality. For
\(\ell_0\in\{1,\ldots,B\}\setminus\mathcal J\), the induction step is
\[
\begin{aligned}
1-\alpha_{\mathrm{aff},\ell_0}
      \prod_{\ell\in\mathcal J}\alpha_{\mathrm{aff},\ell}
&=(1-\alpha_{\mathrm{aff},\ell_0})
  +\alpha_{\mathrm{aff},\ell_0}
     \left(1-\prod_{\ell\in\mathcal J}\alpha_{\mathrm{aff},\ell}\right)\\
&\le (1-\alpha_{\mathrm{aff},\ell_0})
     +1-\prod_{\ell\in\mathcal J}\alpha_{\mathrm{aff},\ell}\\
&\le (1-\alpha_{\mathrm{aff},\ell_0})
     +\sum_{\ell\in\mathcal J}(1-\alpha_{\mathrm{aff},\ell}).
\end{aligned}
\]
Consequently,
\[
\begin{aligned}
\int_{\mathbb R^B}
\left(\sqrt{v_{\boldsymbol s}(y)}
      -\sqrt{v_{\boldsymbol t}(y)}\right)^2\,dy
&\le\sum_{\ell=1}^B\mathcal E_{\mathrm{Hel},\ell}\\
&\le\sum_{\ell=1}^B4\pi^2(s_\ell-t_\ell)^2,
\end{aligned}
\]
where the last inequality is the scalar translation estimate. When \(B=0\), both product densities equal \(1\) on the one-point space \(\mathbb R^0\), and both sides of the bound are zero.
% lean: CausalSmith.Experimentation.ThinnedgraphAdditiveRiskFrontier.hellingerEnergy_pi_le_sum

\item \emph{Total variation for normalized densities.}
Take \(v_+,v_-:\Omega\to[0,\infty)\) satisfying the stated measurability and normalization conditions. For \(\omega_{\mathrm{obs}}\in\Omega\), define
\[
\delta_{\mathrm{dens}}(\omega_{\mathrm{obs}})=v_+(\omega_{\mathrm{obs}})-v_-(\omega_{\mathrm{obs}}),
\qquad
\alpha_{\mathrm{aff}}
=\int_\Omega\sqrt{v_+(\omega_{\mathrm{obs}})v_-(\omega_{\mathrm{obs}})}\,d\mu(\omega_{\mathrm{obs}}).
\]
Normalization makes both densities integrable and gives
\[
\int_\Omega\delta_{\mathrm{dens}}\,d\mu=0.
\]
For any measurable event \(E_{\mathrm{ev}}\subseteq\Omega\), splitting this integral over the event and its complement gives
\[
\int_{E_{\mathrm{ev}}^c}\delta_{\mathrm{dens}}\,d\mu
=-\int_{E_{\mathrm{ev}}}\delta_{\mathrm{dens}}\,d\mu.
\]
Bounding the absolute value of each integral by the integral of the absolute value yields
\[
2\left|\int_{E_{\mathrm{ev}}}\delta_{\mathrm{dens}}\,d\mu\right|
\le
\int_{E_{\mathrm{ev}}}|\delta_{\mathrm{dens}}|\,d\mu
+\int_{E_{\mathrm{ev}}^c}|\delta_{\mathrm{dens}}|\,d\mu
=\int_\Omega|\delta_{\mathrm{dens}}|\,d\mu.
\]
Since the difference of event probabilities is the corresponding set integral, taking the supremum over measurable events proves
\[
\operatorname{TV}(v_+\mu,v_-\mu)
\le\frac12\int_\Omega|v_+-v_-|\,d\mu.
\]

The square roots belong to \(L^2(\mu)\), and pointwise
\[
|v_+-v_-|
=|\sqrt{v_+}-\sqrt{v_-}|\,
 |\sqrt{v_+}+\sqrt{v_-}|.
\]
Cauchy--Schwarz thus gives
\[
\int_\Omega|v_+-v_-|\,d\mu
\le
\left[\int_\Omega(\sqrt{v_+}-\sqrt{v_-})^2\,d\mu\right]^{1/2}
\left[\int_\Omega(\sqrt{v_+}+\sqrt{v_-})^2\,d\mu\right]^{1/2}.
\]
Expanding the first square gives
\[
\int_\Omega(\sqrt{v_+}-\sqrt{v_-})^2\,d\mu
=2(1-\alpha_{\mathrm{aff}}).
\]
For the second square, the pointwise inequality
\[
(\sqrt{v_+}+\sqrt{v_-})^2\le2v_++2v_-
\]
follows from \((\sqrt{v_+}-\sqrt{v_-})^2\ge0\). Therefore,
\[
\int_\Omega(\sqrt{v_+}+\sqrt{v_-})^2\,d\mu\le4.
\]
Combining these bounds and then substituting the energy identity proves
\[
\operatorname{TV}(v_+\mu,v_-\mu)
\le\sqrt{2(1-\alpha_{\mathrm{aff}})}
=
\left[\int_\Omega(\sqrt{v_+}-\sqrt{v_-})^2\,d\mu\right]^{1/2}.
\]
% lean: CausalSmith.Experimentation.ThinnedgraphAdditiveRiskFrontier.density_tv_le_sqrt_hellingerEnergy

\item \emph{Common-marginal averaging.}
Under the common-marginal hypotheses, use marginal and response coordinates
\[
\zeta_{\mathrm{marg}}\in\Xi,
\qquad
\omega_{\mathrm{obs}}\in\Omega,
\]
and define the joint and response laws by
\[
\mathsf P_{\mathrm{joint},\pm}(d\zeta_{\mathrm{marg}},d\omega_{\mathrm{obs}})
=\nu(d\zeta_{\mathrm{marg}})v_\pm(\zeta_{\mathrm{marg}},\omega_{\mathrm{obs}})\mu(d\omega_{\mathrm{obs}}),
\qquad
\mathsf P_{\mathrm{mix},\pm}(d\omega_{\mathrm{obs}})
=\left[\int_\Xi v_\pm(\zeta_{\mathrm{marg}},\omega_{\mathrm{obs}})\,d\nu(\zeta_{\mathrm{marg}})\right]\mu(d\omega_{\mathrm{obs}}).
\]
Joint measurability and integration of nonnegative functions give, for every measurable
\(D_{\mathrm{ev}}\subseteq\Xi\times\Omega\),
\[
\mathsf P_{\mathrm{joint},\pm}(D_{\mathrm{ev}})
=\int_\Xi\int_\Omega
\mathbf1_{D_{\mathrm{ev}}}(\zeta_{\mathrm{marg}},\omega_{\mathrm{obs}})v_\pm(\zeta_{\mathrm{marg}},\omega_{\mathrm{obs}})\,d\mu(\omega_{\mathrm{obs}})\,d\nu(\zeta_{\mathrm{marg}})
=\int_{D_{\mathrm{ev}}}v_\pm(\zeta_{\mathrm{marg}},\omega_{\mathrm{obs}})\,d(\nu\otimes\mu)(\zeta_{\mathrm{marg}},\omega_{\mathrm{obs}}).
\]
Fiber normalization gives
\[
\mathsf P_{\mathrm{joint},\pm}(\Xi\times\Omega)
=\int_\Xi 1\,d\nu=1.
\]
Thus the joint laws are probability laws with integrable densities \(v_\pm\) relative to the sigma-finite product reference measure.

To identify their squared Hellinger distance, define measures on \(\Xi\times\Omega\) by
\[
\Lambda_{\mathrm{joint}}=\nu\otimes\mu,
\qquad
\Lambda_{\mathrm{sum}}
=\mathsf P_{\mathrm{joint},+}+\mathsf P_{\mathrm{joint},-},
\]
and define their Radon--Nikodym density representatives by
\[
r_{\mathrm{sum}}
=\frac{d\Lambda_{\mathrm{sum}}}{d\Lambda_{\mathrm{joint}}},
\qquad
r_\pm
=\frac{d\mathsf P_{\mathrm{joint},\pm}}{d\Lambda_{\mathrm{sum}}}.
\]
These are nonnegative real-valued functions on \(\Xi\times\Omega\), with any infinite derivative values assigned zero; the relevant derivatives are finite almost everywhere relative to their reference measures. The Radon--Nikodym chain rule gives
\[
v_\pm(\zeta_{\mathrm{marg}},\omega_{\mathrm{obs}})=r_\pm(\zeta_{\mathrm{marg}},\omega_{\mathrm{obs}})r_{\mathrm{sum}}(\zeta_{\mathrm{marg}},\omega_{\mathrm{obs}})
\quad\text{for }\Lambda_{\mathrm{joint}}\text{-almost every }(\zeta_{\mathrm{marg}},\omega_{\mathrm{obs}}).
\]
Hence, almost everywhere,
\[
(\sqrt{r_+}-\sqrt{r_-})^2r_{\mathrm{sum}}
=
\left(\sqrt{r_+r_{\mathrm{sum}}}
      -\sqrt{r_-r_{\mathrm{sum}}}\right)^2
=(\sqrt{v_+}-\sqrt{v_-})^2.
\]
Changing reference measure therefore identifies the squared Hellinger distance of the joint laws as
\[
\begin{aligned}
\int_{\Xi\times\Omega}
(\sqrt{r_+}-\sqrt{r_-})^2\,d\Lambda_{\mathrm{sum}}
&=\int_{\Xi\times\Omega}
(\sqrt{v_+}-\sqrt{v_-})^2\,d\Lambda_{\mathrm{joint}}\\
&=\int_\Xi
\left[\int_\Omega
(\sqrt{v_+(\zeta_{\mathrm{marg}},\omega_{\mathrm{obs}})}-\sqrt{v_-(\zeta_{\mathrm{marg}},\omega_{\mathrm{obs}})})^2\,d\mu(\omega_{\mathrm{obs}})\right]d\nu(\zeta_{\mathrm{marg}}).
\end{aligned}
\]
For the last equality, the square-root difference has an integrable square because the joint densities are integrable and their square roots belong to \(L^2(\Lambda_{\mathrm{joint}})\); Fubini applies.

Applying the preceding total-variation bound to these joint densities gives
\[
\operatorname{TV}(\mathsf P_{\mathrm{joint},+},\mathsf P_{\mathrm{joint},-})
\le
\left[
\int_\Xi\int_\Omega
(\sqrt{v_+(\zeta_{\mathrm{marg}},\omega_{\mathrm{obs}})}-\sqrt{v_-(\zeta_{\mathrm{marg}},\omega_{\mathrm{obs}})})^2\,d\mu(\omega_{\mathrm{obs}})\,d\nu(\zeta_{\mathrm{marg}})
\right]^{1/2}.
\]
Finally, for every measurable \(E_{\mathrm{ev}}\subseteq\Omega\), integration of the joint law gives
\[
\mathsf P_{\mathrm{mix},\pm}(E_{\mathrm{ev}})
=\mathsf P_{\mathrm{joint},\pm}(\Xi\times E_{\mathrm{ev}}).
\]
The cylinder event is measurable, so
\[
\begin{aligned}
\left|\mathsf P_{\mathrm{mix},+}(E_{\mathrm{ev}})
      -\mathsf P_{\mathrm{mix},-}(E_{\mathrm{ev}})\right|
&\le
\operatorname{TV}(\mathsf P_{\mathrm{joint},+},\mathsf P_{\mathrm{joint},-}).
\end{aligned}
\]
Taking the supremum over response events proves the asserted mixture bound.
% lean: CausalSmith.Experimentation.ThinnedgraphAdditiveRiskFrontier.common_marginal_hellinger_averaging

\item \emph{Pointwise differentiation away from the endpoints.}
Fix \(w\in\mathbb R\setminus\{-1/4,1/4\}\).
If \(w<-1/4\), then \(\varphi\) agrees with zero throughout a neighborhood of \(w\), and its derivative is \(0=\varphi'(w)\).
For \(w>-1/4\), split further at \(1/4\). If \(w<1/4\), the clamping identity gives
\[
\varphi(w_{\mathrm{loc}})=2\cos(2\pi w_{\mathrm{loc}})
\quad\text{for real }w_{\mathrm{loc}}\text{ in a neighborhood of }w.
\]
The chain rule yields
\[
\left.\frac{d}{dw_{\mathrm{loc}}}\bigl(2\cos(2\pi w_{\mathrm{loc}})\bigr)\right|_{w_{\mathrm{loc}}=w}
=-4\pi\sin(2\pi w)=\varphi'(w).
\]
If \(w>1/4\), then \(\varphi\) again agrees with zero throughout a neighborhood of \(w\), giving derivative \(0=\varphi'(w)\). The two equality cases are excluded by the choice of \(w\).
% lean: CausalSmith.Experimentation.ThinnedgraphAdditiveRiskFrontier.sqrt_cosSqDensity_hasDerivAt

\item \emph{The energy of the derivative representative.}
The endpoint-zero convention in the statement gives the pointwise identity
\[
\varphi'(w)^2
=\mathbf1_{(-1/4,1/4)}(w)
  \bigl(-4\pi\sin(2\pi w)\bigr)^2
\qquad(w\in\mathbb R).
\]
Integrating this indicator, replacing the open interval by the closed interval up to its two null endpoints, and removing the sign inside the square gives
\[
\begin{aligned}
\int_{\mathbb R}\varphi'(w)^2\,dw
&=\int_{(-1/4,1/4)}(-4\pi\sin(2\pi w))^2\,dw\\
&=\int_{[-1/4,1/4]}(4\pi\sin(2\pi w))^2\,dw\\
&=4\pi^2.
\end{aligned}
\]
% lean: CausalSmith.Experimentation.ThinnedgraphAdditiveRiskFrontier.baselineSqrtDerivAt_energy
\end{enumerate}
\end{proof}
